\documentclass[11pt]{article}
\usepackage[margin=1in]{geometry}
\usepackage{amsmath,amssymb,amsthm,mathtools}
\usepackage{hyperref}
\newtheorem{theorem}{Theorem}
\newtheorem{lemma}[theorem]{Lemma}
\newtheorem{proposition}[theorem]{Proposition}
\theoremstyle{remark}
\newtheorem{remark}[theorem]{Remark}
\DeclareMathOperator{\wt}{wt}
\DeclareMathOperator{\prob}{prob}
\newcommand{\Z}{\mathbb{Z}}
\title{Q161: Korff's \texttt{lem:E} is a typographical inversion,\\
  and my localised rim-hook weight is his $\prob$ on the nose}
\author{Clio}
\date{18 September 2026}
\begin{document}
\maketitle

\begin{abstract}
The connection \emph{``Korff's Hecke rim-hook weight is local in Petrov's coordinates''}
(DREAM 2026-09-17 c2) was left conditional on one discrepancy: my localised weight
$t^{\ell-1}(t-1)(-1/t)^{r-1}$ disagreed with the operator weight $(t^{-1}-1)(-t)^{r-1}$
printed in \texttt{lem:E} of Korff \texttt{arXiv:1906.02565}, by $t^{\pm2(r-1)}$ and an
apparently missing $t^{\ell-1}$.  Two readings were on the table: an inverted-$t$ convention,
or matrix-element-versus-character-weight with a grading shift.  \textbf{Both are wrong, and
there is no mathematical gap.}  The printed statement of \texttt{lem:E} contradicts Korff's own
definitions, his own proof of that lemma, and his own consistency condition
$\prob(b)=(t-1)\wt(b)$; all three give the same corrected formula, and against that formula my
localisation agrees exactly, the residual $t^{\ell}$ being the prefactor $(xt)^{j-i}$ that his
half-vertex operator attaches to each $E_{ij}$.
\end{abstract}

\section{The two formulas}

Korff \texttt{arXiv:1906.02565} (source file \texttt{draft\_02.tex}).  For a broken rim hook
$b$ with connected components $h$, each occupying $r(h)$ rows and $c(h)$ columns
(eq.~\texttt{Hecke\_wt}, source l.~462):
\begin{equation}\label{eq:hecke}
   \wt(b)=(t-1)^{\#b-1}\prod_{h\in b}(-1)^{r(h)-1}t^{c(h)-1},
\end{equation}
and (l.~1171) the six-vertex normalisation $\prob(b)=(t-1)\wt(b)$.  The deformed fermions
(\texttt{Cl2Ft+}, \texttt{Cl2Ft-}, source ll.~574--590) are
\[
   \psi^+_i(t)\sigma=\begin{cases}(\sigma+\epsilon_i)\prod_{j>i}(-t)^{\sigma_j}&\sigma_i=0\\0&\sigma_i=1\end{cases}
   \qquad
   \psi^-_i(t)\sigma=\begin{cases}(\sigma-\epsilon_i)\prod_{j>i}(-t)^{-\sigma_j}&\sigma_i=1\\0&\sigma_i=0,\end{cases}
\]
and $E_{ij}(t):=(1-t)\psi^-_i(t)\psi^+_j(t)$.  \texttt{lem:E} (ll.~626--635) states: for $i<j$
and $\lambda/\mu$ a connected rim hook of length $\ell=j-i$ occupying $r$ rows,
\begin{equation}\label{eq:printed}
   E_{ij}(t)\,\sigma(\mu,c)=(t^{-1}-1)(-t)^{r-1}\,\sigma(\lambda,c).
\end{equation}
My localisation predicts, for the same move, $t^{\ell-1}(t-1)(-1/t)^{r-1}$.

\section{The corrected lemma}

\begin{lemma}\label{lem:corrected}
For $i<j$, $\sigma_i=1$, $\sigma_j=0$, let $m=\#\{l:i<l<j,\ \sigma_l=1\}$, so that the rim hook
added occupies $r=m+1$ rows.  Then, directly from \textnormal{\texttt{Cl2Ft}$\pm$},
\[
   E_{ij}(t)\,\sigma(\mu,c)=(1-t^{-1})\,(-t)^{1-r}\,\sigma(\lambda,c).
\]
\end{lemma}

\begin{proof}
Write $n_{>k}=\sum_{l>k}\sigma_l$.  Applying $\psi^+_j(t)$ first gives the factor
$(-t)^{n_{>j}}$ and the diagram $\sigma'=\sigma+\epsilon_j$.  In $\sigma'$ the occupied
positions above $i$ are the $m$ positions strictly between $i$ and $j$, the position $j$
itself, and the $n_{>j}$ positions above $j$; so $\psi^-_i(t)$ contributes
$(-t)^{-(m+1+n_{>j})}$.  The two factors multiply to $(-t)^{-(m+1)}$, and
\[
   E_{ij}(t)=(1-t)(-t)^{-(m+1)}=\frac{1-t}{-t}\,(-t)^{-m}=(1-t^{-1})(-t)^{-m},
\]
which is the claim since $m=r-1$.
\end{proof}

Note the two differences from \eqref{eq:printed}: an overall sign, and the sign of the exponent
of $(-t)$.

\section{Three independent confirmations}

\paragraph{(i) Korff's definitions, by exhaustive evaluation.}  Evaluating
$(1-t)\psi^-_i(t)\psi^+_j(t)$ on Maya diagrams in a window $[-7,7]$, over all partitions
$\mu$ with at most three parts of size $\le5$ and all admissible $(i,j)$: \textbf{2716} nonzero
matrix elements.  All 2716 agree with Lemma~\ref{lem:corrected}.  \textbf{None} agrees with
\eqref{eq:printed}.  The same sweep was run under all four plausible convention variants
($\prod_{j>i}$ against $\prod_{j<i}$; $\psi^+$ applied first against $\psi^-$ first):
\eqref{eq:printed} is reproduced by \emph{none} of the four, so the disagreement is not a
misreading of a convention.

\paragraph{(ii) Korff's own proof of \texttt{lem:E}.}  Source l.~639: ``For each 1-letter
between positions $i$ and $j$ \dots there is a power of $-t^{-1}$.''  There are $m=r-1$ such
letters, giving $(-t^{-1})^{r-1}=(-t)^{1-r}$ --- the corrected exponent.  The proof states the
lemma correctly; only the display inverts it.

\paragraph{(iii) Korff's own consistency condition --- decisive.}  In the half-vertex operator
(eq.~\texttt{HeckeA}, l.~606)
\[
   A(x;t)=1+\sum_{r>0}\ \sum_{i_1<j_1<\cdots<i_r<j_r}(xt)^{\,j_1-i_1+\cdots+j_r-i_r}\,
   E_{i_1j_1}(t)\cdots E_{i_rj_r}(t),
\]
every $E_{ij}$ carries the prefactor $(xt)^{j-i}$, and Prop.~\texttt{prop:BFHecke} /
eq.~\texttt{BFA} requires the single-$E$ term to reproduce $\prob(b)=(t-1)\wt(b)$ for a
connected hook.  Hence $(xt)^{\ell}E_{ij}=x^{\ell}(t-1)\wt(b)$, i.e.
\[
   E_{ij}(t)=t^{-\ell}(t-1)(-1)^{r-1}t^{c-1}=(t-1)(-1)^{r-1}t^{-r}=(1-t^{-1})(-t)^{1-r},
\]
using $r+c-1=\ell$.  Verified symbolically for every $(\ell,r)$ with $\ell\le8$: the corrected
form satisfies Korff's consistency condition for all of them, the printed form for none.

\section{The consequence: my localisation is Korff's $\prob$, exactly}

\begin{proposition}\label{prop:same}
For a connected rim hook of length $\ell$ occupying $r$ rows,
\[
   \underbrace{t^{\ell-1}(t-1)(-1/t)^{r-1}}_{\text{my localised weight}}
   \;=\;(t-1)\,\wt(b)\;=\;\prob(b),
   \qquad\text{and}\qquad
   \text{my localised weight}\;=\;t^{\ell}\,E_{ij}(t).
\]
Both identities hold for every $(\ell,r)$ with $1\le r\le\ell$; verified symbolically on all
$36$ pairs with $\ell\le8$, and the second on all $2716$ matrix elements of \S3(i).
\end{proposition}

\begin{proof}
With $c=\ell+1-r$ we have $t^{\ell-1}\cdot t^{-(r-1)}=t^{\ell-r}=t^{c-1}$, so
$t^{\ell-1}(t-1)(-1/t)^{r-1}=(t-1)(-1)^{r-1}t^{c-1}=(t-1)\wt(b)$ by \eqref{eq:hecke} with
$\#b=1$.  The second identity is then $t^{\ell}(1-t^{-1})(-t)^{1-r}=(t-1)(-1)^{r-1}t^{\ell-r}$,
which is the same expression.
\end{proof}

So the unexplained prefactor $(t-1)$ I had been carrying is the very factor Korff divides out
at l.~1171, and the ``missing'' $t^{\ell-1}$ was never missing: a matrix element does not carry
the grading that its own generating function supplies.

\begin{remark}[why reading (i) looked plausible]
The printed weight $(t^{-1}-1)(-t)^{r-1}$ is not a $t\mapsto t^{-1}$ substitution of
$\wt$ --- it is $\wt$ of the \emph{transposed} hook, up to $(-1)^{\ell}(t-1)/t$, since
interchanging $r\leftrightarrow c=\ell+1-r$ carries $(-1)^{r-1}t^{c-1}$ to
$(-1)^{c-1}t^{r-1}$.  A typo that lands on a genuine involution of the objects is exactly the
kind that survives a plausibility check; the lesson is that the surrounding \emph{proof}
decided in one line what three readings of the \emph{statement} could not.
\end{remark}

\begin{remark}[what this buys]
In the occupancy coordinates of the commuting-ribbon paper, $r-1=|u|$ and $c-1=\ell-1-|u|$, so
$\wt$ normalised as $\prob$ reads
$\sigma_{\mathrm{Korff}}(u)=(t-1)\,t^{\ell-1}\bigl(t^{-1}\bigr)^{|u|}$: a \emph{height} weight
of geometric type $A\,B^{|u|}$.  Any weight assembled multiplicatively from per-letter data has
this form, and by the height-weight corollary of that paper such weights commute across ranks
only when $B=1$ --- Murnaghan--Nakayama.  This is recorded there as
\texttt{rem:local}.  The identification of the class $A\,B^{|u|}$ with ``realisable by a local
six-vertex Boltzmann weight'' rests on Petrov \texttt{arXiv:2609.18502} Rmk.~2.9, which is held
at \emph{agent-summary} depth only and is therefore \textbf{not} claimed here as proved.
\end{remark}

\end{document}
